% Generated draft: requires mathematical and source review.
\documentclass[11pt]{article}
\usepackage[margin=25mm]{geometry}
\usepackage{fontspec}
\setmainfont{Latin Modern Roman}
\usepackage{amsmath,amssymb,mathtools,unicode-math}
\setmathfont{Latin Modern Math}
\usepackage{enumitem,needspace,xurl,hyperref}
\setlength{\parindent}{0pt}
\setlength{\parskip}{5pt}
\setlength{\emergencystretch}{3em}
\newcommand{\N}{\mathbb{N}}
\newcommand{\Z}{\mathbb{Z}}
\newcommand{\Q}{\mathbb{Q}}
\newcommand{\R}{\mathbb{R}}
\newcommand{\C}{\mathbb{C}}
\newcommand{\F}{\mathbb{F}}
\newcommand{\A}{\mathbb{A}}
\newcommand{\PP}{\mathbb{P}}
\newcommand{\E}{\mathbb{E}}
\newcommand{\eps}{\varepsilon}
\newcommand{\dd}{\,\mathrm{d}}
\begin{document}
\section*{UnsolvedMath ID: 30002535}
\textbf{Source problem code:} catalogue entry
\subsection{Definitions and mathematical statement}
A terminal projective threefold is a normal projective variety \(Y\) of dimension \(3\), with \(K_Y\) \(\Q\)-Cartier and with strictly positive discrepancies for all exceptional divisorial valuations over \(Y\). The supplied record does not specify the ground field or its characteristic, or whether \(Y\) is assumed to be \(\Q\)-factorial.

Fix ambient hypotheses and conventions under which the \(K_Y\)-minimal model program (MMP) is defined. A reachable minimal-model outcome is a birational map \(\phi\colon Y\dashrightarrow Y'\) obtained by a finite sequence, possibly empty, of permitted \(K\)-negative divisorial contractions and flips, whose terminal output \(Y'\) is a minimal model; in particular, \(K_{Y'}\) is nef, meaning \(K_{Y'}\cdot C\geq 0\) for every integral curve \(C\subseteq Y'\). Any additional factoriality and singularity requirements on the output must be specified by the chosen MMP conventions. A Mori fibre space is not counted as a minimal-model outcome here.

Let \(\sim\) be a specified equivalence relation on such outcomes, and write
\[
\mathcal M_{\sim}(Y)
 =\bigl\{(Y',\phi): (Y',\phi)\text{ is a reachable minimal-model outcome of }Y\bigr\}/\!\sim.
\]
For an ample divisor \(H\) on \(Y\), let \(\mathcal M_{\sim}(Y;H)\subseteq\mathcal M_{\sim}(Y)\) consist of the outcomes obtainable by an MMP with scaling of \(H\), under conventions for which this scaling program is defined.

The supplied conjecture asks whether
\[
\#\mathcal M_{\sim}(Y)<\infty
\qquad\text{for every terminal projective threefold }Y.
\]
It also asks, in particular, whether
\[
\#\mathcal M_{\sim}(Y;H)<\infty
\qquad\text{for every such }Y\text{ and every fixed ample divisor }H\text{ on }Y.
\]
These are finiteness questions about outcomes, not assertions that every such \(Y\) has a minimal model. If no minimal-model outcome is reachable, the corresponding set is empty.

\emph{Unresolved qualifications.} The record does not determine the equivalence relation \(\sim\): counting varieties up to isomorphism can differ from counting models together with their birational markings from \(Y\). It also leaves the ambient field, factoriality conventions, permissible initial modifications, and precise scaling hypotheses unspecified. Thus the displayed formulation is conditional on those choices and does not supply a uniquely determined standalone conjecture. The fixed-\(H\) reading of the scaling clause is displayed explicitly; the record does not clarify whether that clause instead intends to count outcomes jointly over varying ample divisors.
\subsection{Short English statement}
The supplied conjecture asks whether a terminal projective threefold has only finitely many minimal models obtainable by the canonical-divisor MMP, including the version with scaling of an ample divisor. The ground field, factoriality assumptions, and meaning of counting distinct models require clarification. The supplied open status is dated curation, not an independently verified conclusion.
\subsection{Sources}
\begin{itemize}
\item Local source registry: \texttt{data/unsolvedmath/problems.json}, with \texttt{catalog/problems.json} as the supplementary registry; supplied record for permanent UnsolvedMath ID \(30002535\), exact source problem code \texttt{catalogue entry}, title \emph{Finiteness of Minimal Models of Terminal Threefolds}. These registry references identify the supplied record and have not been independently inspected.
\item Supplied attribution: \emph{Mini-Workshop: Negative Curves on Algebraic Surfaces} (2014), \emph{Oberwolfach Reports}, pp.~555--593, DOI \texttt{10.4171/owr/2014/10}. The record attributes the problem to the workshop-report collection and gives proposed year 2014; no individual proposer is supplied. Source link: \url{https://doi.org/10.4171/owr/2014/10}.
\item Supplied literature source, labelled ``Oberwolfach Report 10, 2014'': \url{https://doi.org/10.4171/owr/2014/10}.
\item Supplied literature source, labelled ``Oberwolfach Report 10 PDF'': \url{https://publications.mfo.de/bitstream/handle/mfo/3399/OWR_2014_10.pdf?isAllowed=y\&sequence=1}.
\item Literature assessment dated exactly \texttt{2026-08-21}: the supplied curation labels the problem open and states that finiteness of minimal-model outcomes for terminal projective threefolds, including the ample-scaling version, is not settled in general. This is supplied dated curation derived from the OWR record, not an independently verified literature assessment. The record expressly requests independent verification of time-sensitive status and cited claims.
\end{itemize}
\end{document}
